% source: manti.01.ma, 1ère SM, série suites (Drive 0B6pfSJcEVng0c2x5Unp5VU5pR0k), p.1-3
\ex{Suites}{c2x5}{Ex. 1}{8}
Soit $(U_n)_{n\in\mathbb{N}^*}$ la suite définie par $U_n=\displaystyle\sum_{k=1}^{n}\frac{1}{k^2}$.
Montrer que pour tout entier $k\geq 2$ : $\dfrac{1}{k^2}<\dfrac{1}{k-1}-\dfrac{1}{k}$, puis montrer que $(\forall n\in\mathbb{N}^*)\ U_n<2$.
\fin
\ex{Suites}{c2x5}{Ex. 2}{7}
Soit $(U_n)_{n\geq 2}$ la suite définie par $U_n=\displaystyle\sum_{k=2}^{n}\frac{1}{k^2-1}$.
\begin{enumerate}
\item Calculer $U_2$ et $U_3$.
\item Étudier la monotonie de la suite $(U_n)_{n\geq2}$.
\item Vérifier que $\dfrac{1}{k^2-1}=\dfrac12\left(\dfrac{1}{k-1}-\dfrac{1}{k+1}\right)$, puis en déduire que $(U_n)_{n\geq2}$ est majorée.
\end{enumerate}
\fin
\ex{Suites}{c2x5}{Ex. 3}{9}
On considère la suite $(U_n)_{n\geq1}$ définie par $U_n=\dfrac{1}{\sqrt n}\displaystyle\sum_{k=1}^{n}\frac{1}{\sqrt k}$.
\begin{enumerate}
\item Calculer $U_1$ et $U_2$, et montrer que $U_2>U_1$.
\item \begin{enumerate}
\item Montrer que pour tout entier naturel non nul $k$ : $\dfrac{1}{\sqrt{k+1}}\leq 2\left(\sqrt{k+1}-\sqrt k\right)$.
\item En déduire que $(U_n)_{n\geq1}$ est majorée par $2$.
\end{enumerate}
\end{enumerate}
\fin
\ex{Suites}{c2x5}{Ex. 4}{8}
On considère la suite $(U_n)_{n\in\mathbb{N}}$ définie par $U_0=3$ et $U_{n+1}=3-\dfrac{9}{4U_n}$.
\begin{enumerate}
\item Montrer que $(\forall n\in\mathbb{N})\ U_n>\dfrac32$.
\item Étudier la monotonie de la suite $(U_n)_{n\geq0}$.
\item On pose $V_n=\dfrac{2}{2U_n-3}$.
\begin{enumerate}
\item Montrer que $(V_n)_{n\geq0}$ est une suite arithmétique.
\item Montrer que $U_n=\dfrac32\left(\dfrac{n+2}{n+1}\right)$.
\end{enumerate}
\end{enumerate}
\fin
\ex{Suites}{c2x5}{Ex. 5}{8}
On considère la suite $(U_n)_{n\in\mathbb{N}}$ définie par $U_0=\dfrac43$ et $U_{n+1}=\dfrac{6U_n-3}{2U_n+1}$.
\begin{enumerate}
\item Montrer que $(\forall n\in\mathbb{N})\ 1<U_n<\dfrac32$ \quad (remarquer que $U_{n+1}=3-\dfrac{6}{2U_n+1}$).
\item Étudier la monotonie de la suite $(U_n)_{n\in\mathbb{N}}$.
\item On pose $V_n=\dfrac{2U_n-3}{U_n-1}$. Montrer que $(V_n)_{n\geq0}$ est une suite géométrique, puis calculer $U_n$ en fonction de $n$.
\end{enumerate}
\fin
\ex{Suites}{c2x5}{Ex. 6}{10}
On considère la suite $(U_n)_{n\in\mathbb{N}}$ définie par $U_0=\dfrac12$ et $U_{n+1}=\dfrac{2U_n^2}{1+U_n^3}$.
\begin{enumerate}
\item Montrer que $(\forall n\in\mathbb{N})\ 0<U_n\leq\dfrac12$.
\item Étudier la monotonie de la suite $(U_n)_{n\in\mathbb{N}}$.
\item Montrer que $U_{n+1}\leq\dfrac89U_n$, puis en déduire que $U_n\leq\dfrac12\left(\dfrac89\right)^n$.
\end{enumerate}
\fin
\ex{Suites}{c2x5}{Ex. 7}{11}
Soit $(U_n)_{n\geq1}$ la suite définie par $U_1=1$ et $U_{n+1}=\dfrac{4n-2}{n+1}+\dfrac{3n}{n+1}U_n$.
\begin{enumerate}
\item Montrer que $(\forall n\in\mathbb{N}^*)\ U_n>-2$.
\item Étudier la monotonie de la suite $(U_n)_{n\geq1}$.
\item On pose $V_n=n\,(2+U_n)$.
\begin{enumerate}
\item Montrer que $(V_n)_n$ est une suite géométrique en précisant sa raison, puis calculer $U_n$ en fonction de $n$.
\item Calculer la somme $S_n=\displaystyle\sum_{k=1}^{n}kU_k$.
\end{enumerate}
\end{enumerate}
\fin
\ex{Suites}{c2x5}{Ex. 8}{10}
Soit $(U_n)_{n\in\mathbb{N}}$ la suite définie par $U_0=3$, $U_1=3$ et $U_{n+2}=\dfrac43U_{n+1}-\dfrac49U_n$.
\begin{enumerate}
\item Calculer $U_2$ et $U_3$.
\item On pose $V_n=U_{n+1}-\dfrac23U_n$. Montrer que $(V_n)_n$ est une suite géométrique, préciser sa raison et calculer $V_n$ en fonction de $n$.
\item En déduire que $U_{n+1}=\dfrac23U_n+\left(\dfrac23\right)^n$.
\item On pose $W_n=\left(\dfrac32\right)^nU_n$ \ (l'original porte $V_n$ [?]).
\begin{enumerate}
\item Montrer que $(W_n)_{n\in\mathbb{N}}$ est une suite arithmétique en précisant sa raison $r$.
\item En déduire que $U_n=\left(\dfrac23\right)^{n-1}(n+2)$.
\end{enumerate}
\end{enumerate}
\fin
\ex{Suites}{c2x5}{Ex. 9}{8}
\begin{enumerate}
\item Soit $(U_n)_{n\in\mathbb{N}}$ une suite telle que, pour tout $n\in\mathbb{N}$ : $S_n=U_0+U_1+\dots+U_{n-1}=\dfrac{2n^2+n}{2}$.
Calculer $U_0$ et $U_1$, puis montrer que $(U_n)$ est une suite arithmétique et donner $U_n$ en fonction de $n$.
\item Soit $(U_n)_{n\in\mathbb{N}}$ une suite arithmétique de raison $r\neq0$ telle que $U_0=4$ et que $U_1$, $U_4$, $U_{12}$ soient des termes consécutifs d'une suite géométrique.
Montrer que $r=5$, puis calculer $U_0+U_1+\dots+U_{12}$.
\item Soit $(U_n)_{n\in\mathbb{N}}$ une suite arithmétique de raison $r>0$ telle que $U_0+U_1+U_2=15$ et $U_0^2+U_1^2+U_2^2=107$.
Montrer que $U_1=5$ et en déduire que $r=4$, puis calculer $U_0+U_1+\dots+U_n$.
\item Soit $(V_n)_n$ une suite géométrique telle que $V_0V_1V_2=27$ et $V_0+V_1+V_2=13$.
Montrer que $V_1=3$ et en déduire que $q=3$ ou $q=\dfrac13$. Calculer, dans chaque cas, $V_0+V_1+\dots+V_{n-1}$.
\end{enumerate}
\fin
\ex{Suites}{c2x5}{Ex. 10}{9}
Soit $(U_n)_{n\in\mathbb{N}}$ la suite définie par $U_0=1$ et $U_{n+1}=2U_n-\dfrac12n+\dfrac12$. On pose $V_n=2U_n-n$.
\begin{enumerate}
\item Calculer $U_1$ et montrer par récurrence que $(\forall n\in\mathbb{N})\ U_n\geq\dfrac n2$.
\item Étudier la monotonie de la suite $(U_n)_{n\in\mathbb{N}}$.
\item Montrer que $(V_n)_n$ est une suite géométrique en précisant sa raison.
\item \begin{enumerate}
\item Montrer que $U_n=2^n+\dfrac12n$.
\item Montrer que $U_0+U_1+\dots+U_{n-1}=2^n-1+\dfrac{n(n-1)}{4}$.
\end{enumerate}
\end{enumerate}
\fin
% Ex. 11, 12, 13 : identiques à la série française suites_s4 (Ex. 7, 8, 9) -> non repris
\ex{Suites}{c2x5}{Ex. 14}{11}
Soit $(U_n)_n$ la suite définie par $U_0=\dfrac12$ et $U_{n+1}=\sqrt{1+\dfrac32U_n}$.
\begin{enumerate}
\item Calculer $U_1$ et montrer que $(\forall n\in\mathbb{N})\ 0<U_n<2$.
\item Étudier la monotonie de la suite $(U_n)_n$.
\item \begin{enumerate}
\item Montrer que $(\forall n\in\mathbb{N})\ |U_{n+1}-2|\leq\dfrac35|U_n-2|$.
\item En déduire que $(\forall n\in\mathbb{N})\ |U_n-2|\leq\dfrac32\left(\dfrac35\right)^n$.
\end{enumerate}
\end{enumerate}
\fin
